\noindent \textbf{\large 1. Conexão e Alimentação} \\
Para iniciar, conecte a placa ESP32 ao seu computador utilizando um cabo Micro-USB de boa qualidade.

\begin{itemize}
    \item \textbf{Reconhecimento:} O computador deverá emitir um som de novo dispositivo conectado.
    \item \textbf{Driver:} Caso seja a primeira vez, certifique-se de ter os drivers instalados (geralmente CP210x ou CH340).
    \item \textbf{Alimentação:} Uma vez conectada, o LED de power da placa acenderá, indicando que ela está recebendo energia.
\end{itemize}

\noindent \textbf{\large 2. Preparação do Ambiente (PlatformIO)} \\
Para enviar o código para a placa, você precisará preparar o software:

\begin{itemize}
    \item \textbf{Placa:} Vá em \textit{Ferramentas > Placa} e selecione \textbf{DOIT ESP32 DEVKIT V1}.
    \item \textbf{Porta:} Vá em \textit{Ferramentas > Porta} e selecione a porta COM que apareceu após conectar o cabo.
    \item \textbf{Bibliotecas:} Instale as bibliotecas \textbf{DHT sensor library (Adafruit), WiFi, WebServer e SPIFFS} através do Gerenciador de Bibliotecas.
\end{itemize}

\noindent \textbf{\large 3. Upload do Código e Arquivos} \\
Este sistema utiliza dois tipos de memória: a do código e a de arquivos (SPIFFS).

\begin{itemize}
    \item \textbf{Código (.ino):} Clique na seta de "Upload"  na IDE para gravar a lógica.
    \item \textbf{Interface Web (pasta data):} Use a ferramenta \textbf{ESP32 Sketch Data Upload} para enviar os arquivos index.html, css e js.
    \item \textit{Nota: Se não enviar os arquivos, a página web aparecerá em branco ou com erro 404.}
\end{itemize}

\noindent \textbf{\large 4. Inicialização e Conexão Wi-Fi} \\
Após o upload, abra o Monitor Serial (lupa no canto superior direito) e configure a velocidade para \textbf{115200 baud}.

\begin{itemize}
    \item A placa tentará se conectar ao Wi-Fi configurado no código.
    \item Quando conectada, ela exibirá uma mensagem: \textbf{IP: 192.168.x.x}. Anote este número.
\end{itemize}

\noindent \textbf{\large 5. Acesso ao Dashboard} \\
\begin{enumerate}
    \item Abra o navegador (Chrome, Edge ou Firefox) em um dispositivo conectado à mesma rede Wi-Fi.
    \item Digite o endereço IP anotado na barra de endereços e pressione Enter.
    \item A interface do monitoramento da \textbf{EletronJun} aparecerá na tela.
\end{enumerate}

\noindent \textbf{\large 6. Operação do Sistema} \\
Agora você pode interagir com o sistema através da página:

\begin{itemize}
    \item \textbf{Troca de Modos:} Selecione entre Desligado, Automático ou Ligado. No modo \textbf{Automático}, o sistema salva as leituras a cada 2 minutos no arquivo interno.
    \item \textbf{Configuração:} No modo \textbf{Ligado}, você pode alterar a temperatura máxima permitida e mudar a unidade de Celsius para Fahrenheit. Clique em \textbf{"Salvar Configuração"} para aplicar.
    \item \textbf{Alertas:} Se a temperatura ultrapassar o limite ou a umidade cair abaixo de 30\%, os LEDs na placa piscarão rapidamente e a interface mostrará um aviso visual.
    \item \textbf{Histórico (Logs):} No final da página, clique em \textbf{Baixar Histórico de Logs} para baixar o arquivo \textit{log.txt} com todas as medições salvas.
\end{itemize}

\vspace{1cm}

\noindent \textbf{\large Resumo Técnico dos Componentes} \\
\begin{itemize}
    \item \textbf{ESP32:} O cérebro que processa os dados e hospeda o site.
    \item \textbf{DHT11:} O sensor que mede o ambiente.
    \item \textbf{LEDs:} Indicadores físicos de modo e perigo.
    \item \textbf{SPIFFS:} Memória interna onde o site e o histórico de texto ficam guardados.
\end{itemize}